\documentclass[12pt]{article}

% --- Packages ---
\usepackage[a4paper, margin=1in]{geometry}
\usepackage{setspace}
\singlespacing
\usepackage{amsmath, amssymb}
\usepackage{graphicx}
\usepackage{booktabs}
\usepackage{tabularx}
\usepackage{float}
\usepackage{hyperref}
\usepackage{cite}
\usepackage{enumitem}
\usepackage[T1]{fontenc}
\usepackage{times}
\usepackage{CJKutf8}

% --- Title ---
\title{\textbf{Dynamic EDCA Parameter Adaptation for\\
Starvation Avoidance in IEEE 802.11 WLANs Proposal}}
\author{
  M143140017 \begin{CJK}{UTF8}{bsmi}王承煜\end{CJK}\\
  M143040081 \begin{CJK}{UTF8}{bsmi}林憲逸\end{CJK}\\[1em]
  CSE625-Wireless Communication Network:\\
  Construction and Performance Simulation\\
  }
\date{}

\begin{document}
\maketitle

% ============================================================
\section{Background}
% ============================================================

The Enhanced Distributed Channel Access (EDCA) mechanism in
IEEE 802.11 provides quality-of-service (QoS) differentiation
by classifying traffic into four access categories (ACs): Voice
(AC\_VO), Video (AC\_VI), Best Effort (AC\_BE), and Background
(AC\_BK)~\cite{ieee80211}. Each AC is assigned distinct channel
access parameters---Arbitration Inter-Frame Space Number (AIFSN),
contention window (CWmin/CWmax), and Transmission Opportunity
(TXOP) limit---that determine the priority with which a station
competes for the wireless medium.

However, the static nature of these parameters leads to a
well-documented problem: under moderate-to-heavy load,
high-priority ACs monopolize channel access, causing low-priority
ACs to experience severe \emph{starvation} with near-zero
throughput and excessive packet loss. Ugwu et
al.~\cite{ugwu2022} confirmed through MATLAB simulation that
AIFSN has the greatest influence on this disparity. Mammeri et
al.~\cite{mammeri2021} proposed a Starvation avoidance-based
Dynamic Multichannel Access (SDMA) method for 802.11ac, but it
relies on channel bonding and does not apply to single-channel
scenarios.

EDCA remains the foundational channel access mechanism in
Wi-Fi~6 (802.11ax)~\cite{ieee80211ax} and Wi-Fi~7
(802.11be)~\cite{yi2025}. A significant portion of recent
research emphasizes OBSS/inter-BSS interference management
(BSS coloring, OBSS-PD, spatial
reuse)~\cite{tuan2023} and multi-link EDCA optimization under
MLO. While these efforts address system-level resource
allocation across BSSs or links, lightweight MAC-layer
mechanisms targeting intra-BSS AC starvation within a single
BSS remain comparatively underexplored.

Recent deep reinforcement learning (DRL) approaches show
promising results: Zuo et al.~\cite{zuo2025} proposed PDCF-DRL,
maintaining normalized throughput above 76\% and collision rates
below 18\% across 20--120 stations, compared to standard EDCA's
collision rates of 64\%--85\% and throughput of only 13\%--32\%.
Du et al.~\cite{du2024} combined federated learning with DDPG,
reducing MAC delay by up to 45.9\%. Li et al.~\cite{li2024}
proposed ReinWiFi for application-layer QoS optimization.
However, DRL methods require high computational resources and
lengthy training convergence, while application-layer solutions
bypass the MAC layer and do not address the root cause.

% ============================================================
\section{Problem Definition}
% ============================================================

We define starvation for a low-priority access category $AC_i$
($i \in \{BE, BK\}$) as:
\begin{equation}
  \text{Starvation}(AC_i) \iff
  \left( \frac{Q_i}{Q_{\text{cap}}} > Q_{th} \right)
  \;\lor\;
  \left( P_{\text{loss},i} > P_{th} \right)
\end{equation}
where $Q_i / Q_{\text{cap}}$ is the queue occupancy ratio and
$P_{\text{loss},i}$ is the packet loss rate within a monitoring
interval. The two metrics are combined with logical OR because
they capture complementary symptoms: queue buildup reflects
sustained channel deprivation, while packet loss captures
starvation manifesting as frame drops.

The standard EDCA assigns static parameters (e.g., AC\_VO
$CWmin = 3$ vs.\ AC\_BE $CWmin = 15$), giving high-priority
ACs a permanent structural advantage that compounds under heavy
load~\cite{ugwu2022}. Table~\ref{tab:limitations} summarizes
the limitations of existing approaches.

\begin{table}[H]
\centering
\caption{Limitations of existing approaches.}
\label{tab:limitations}
\small
\begin{tabularx}{\textwidth}{l l X}
\toprule
\textbf{Approach} & \textbf{Representative}
  & \textbf{Key Limitation} \\
\midrule
Dynamic multichannel
  & Mammeri~\cite{mammeri2021}
  & Relies on 802.11ac channel bonding; not applicable
    to single-channel BSS \\
OBSS QoS improvement
  & Tuan~\cite{tuan2023}
  & Targets inter-BSS interference; does not address
    intra-BSS AC starvation \\
DRL-based schemes
  & Zuo~\cite{zuo2025}, Du~\cite{du2024}
  & High computational complexity; slow training
    convergence \\
Application-layer
  & Li~\cite{li2024}
  & Operates above MAC; does not adjust EDCA
    parameters directly \\
\bottomrule
\end{tabularx}
\end{table}

Existing dynamic EDCA approaches typically target different
network types (e.g., vehicular or IoT), adjust only one or two
parameters (usually CW alone), or lack a formal starvation
detection criterion. There remains no lightweight, MAC-layer
mechanism that can sense queue state in real time and dynamically
adjust EDCA parameters for starvation avoidance in infrastructure
WLANs. QAD-EDCA addresses these gaps by combining a
dual-indicator starvation detection formula (queue occupancy and
packet loss rate) with simultaneous three-dimensional parameter
adjustment (AIFSN, CWmin, TXOP), complemented by an exponential
recovery mechanism and a QoS safety valve, forming an integrated
closed-loop starvation avoidance system.

\subsection{Proposed Approach}

We propose \textbf{Queue-Aware Dynamic EDCA (QAD-EDCA)}, a
lightweight, rule-based adaptation mechanism running on the
access point (AP). QAD-EDCA operates as a centralized monitor
with $O(1)$ computational complexity per cycle---it reads
4~AC queue lengths and loss rates, evaluates threshold
conditions, and applies deterministic adjustments every
monitoring interval ($T_{\text{mon}} = 100$\,ms), synchronized
with the standard beacon interval so that no extra overhead is
introduced.

\paragraph{Three-Stage Adjustment.}
Upon detecting starvation, QAD-EDCA simultaneously applies
three complementary strategies targeting different phases of
the EDCA channel access process:
\begin{enumerate}[noitemsep]
  \item \textbf{Reduce AIFSN for low-priority ACs.}
    $AIFSN_i = \max(AIFSN_i - \Delta_{\text{AIFS}},\; 2)$
    for $i \in \{BE, BK\}$. This shortens pre-contention wait
    time. The lower bound of~2 matches AC\_VO/AC\_VI defaults,
    ensuring low-priority ACs never become more aggressive than
    high-priority ones.
  \item \textbf{Increase CWmin for high-priority ACs.}
    $CWmin_i = \min(CWmin_i \times \alpha_{CW},\;
    CWmin_{\text{default},BE})$ for $i \in \{VO, VI\}$. This
    broadens the backoff window, reducing high-priority
    dominance. When CWmin increases from 3 to 15, the per-slot
    transmission probability drops from 50\% to 12.5\%.
  \item \textbf{Reduce TXOP for high-priority ACs.}
    $TXOP_i = \max(TXOP_i \times \beta_{\text{TXOP}},\;
    TXOP_{\min})$ for $i \in \{VO, VI\}$. This forces
    high-priority STAs to release the channel sooner, creating
    more access opportunities for low-priority traffic.
\end{enumerate}
The three strategies act on different dimensions---AIFSN
controls \emph{when} contention begins, CWmin controls
\emph{how long} the backoff lasts, and TXOP controls
\emph{how long} the channel is held---so applying them
concurrently provides the fastest starvation response.

\paragraph{Recovery Mechanism.}
When the starvation condition is no longer met, parameters
do not snap back to their defaults instantly. Instead, they
recover gradually via exponential decay:
$P(t{+}1) = P(t) + \gamma \cdot (P_{\text{default}} - P(t))$
with $\gamma = 0.3$, meaning each monitoring cycle closes 30\%
of the remaining gap between the current value and the default.
After 5~cycles ($\sim$500\,ms), parameters reach approximately
83\% of their default values ($1 - 0.7^5 \approx 0.83$).
Gradual recovery is essential because an instantaneous reset
would immediately restore high-priority dominance, re-triggering
starvation and creating a ping-pong oscillation between
adjustment and recovery states.

\paragraph{QoS Safety Valve.}
After each adjustment, the algorithm checks whether
high-priority delay constraints are still satisfied. If
$d_{VO} > 150$\,ms or $d_{VI} > 300$\,ms, it partially rolls
back: $P_{\text{adj}} \leftarrow (P_{\text{adj}} +
P_{\text{default}}) / 2$. This ensures QoS preservation always
takes precedence over starvation mitigation.

\paragraph{Evaluation Plan.}
The scheme will be evaluated in \textbf{OMNeT++ 6.1} with the
\textbf{INET Framework 4.5}, comparing against standard EDCA,
tuned static EDCA~\cite{ugwu2022}, and
PDCF-DRL~\cite{zuo2025}. Key metrics include per-AC throughput,
delay, packet loss rate, and Jain's fairness
index~\cite{jain1984}.

\paragraph{Positioning.}
It is important to note that QAD-EDCA is \textbf{not} designed
to outperform DRL in raw throughput optimization. Rather, it
fills the gap between static EDCA and computationally expensive
DRL approaches by offering a practical, immediately deployable
alternative. DRL-based schemes such as
PDCF-DRL~\cite{zuo2025} achieve superior throughput through
learned policies, but require extensive offline training,
significant computational resources, and retraining when network
conditions change. QAD-EDCA, as a zero-training,
$O(1)$-complexity rule engine, can be deployed on any standard
AP without hardware upgrades. Our comparison with DRL therefore
aims to characterize the \emph{performance--complexity
trade-off}, not to claim superiority in absolute throughput.

% ============================================================
\section{Project Plan}
% ============================================================

\subsection{Timeline}

\begin{table}[H]
\centering
\small
\renewcommand{\arraystretch}{1.2}
\begin{tabularx}{\textwidth}{c l X}
\toprule
\textbf{Weeks} & \textbf{Period} & \textbf{Tasks} \\
\midrule
1--2 & 4/15 -- 4/28
  & Environment setup; implement baseline EDCA simulation;
    reproduce starvation phenomenon under heavy load \\
3--4 & 4/29 -- 5/12
  & Implement QAD-EDCA algorithm in OMNeT++/INET;
    integrate with Edcaf module; unit testing \\
5--6 & 5/13 -- 5/26
  & Execute full experiment matrix (varying STA count,
    load level, threshold parameters); implement
    comparison schemes \\
7--8 & 5/27 -- 6/10
  & Analyze results and generate figures; write final
    report; prepare presentation slides \\
\bottomrule
\end{tabularx}
\end{table}

\subsection{Team Responsibilities}

All members share responsibilities across roles.

\vspace{0.5em}
\noindent
\begin{tabular}{@{} p{4cm} p{11.5cm} @{}}
\toprule
\textbf{Role} & \textbf{Responsibilities} \\
\midrule
Algorithm and Implementation
  & QAD-EDCA algorithm design;
    C++ implementation in OMNeT++/INET \\[0.3em]
Simulation
  & Network topology setup; traffic configuration;
    running experiments \\[0.3em]
Analysis
  & Result parsing; statistical analysis;
    figure generation \\[0.3em]
Report
  & Literature survey; proposal and final report
    writing; presentation \\
\bottomrule
\end{tabular}

% ============================================================

\emergencystretch=3em
\begin{thebibliography}{10}

\bibitem{ieee80211}
IEEE Std 802.11-2020, ``Wireless LAN Medium Access Control (MAC) and Physical Layer (PHY) Specifications,'' IEEE, 2020.

\bibitem{ugwu2022}
G.~O.~Ugwu, U.~N.~Nwawelu, and M.~A.~Ahaneku, ``Effect of service differentiation on QoS in IEEE 802.11e enhanced distributed channel access: a simulation approach,'' \emph{J.~Eng.~Appl.~Sci.}, vol.~69, no.~1, pp.~1--15, 2022.

\bibitem{mammeri2021}
S.~Mammeri, M.~Yazid, R.~Kacimi, and L.~Bouallouche-Medjkoune, ``Starvation avoidance-based dynamic multichannel access for low priority traffics in 802.11ac communication systems,'' \emph{Comput.~Electr.~Eng.}, vol.~90, art.~106942, 2021.

\bibitem{tuan2023}
Y.~P.~Tuan, L.~A.~Chen, T.~Y.~Lin, et~al., ``Improving QoS mechanisms for IEEE 802.11ax with overlapping basic service sets,'' \emph{Wireless Netw.}, vol.~29, pp.~387--401, 2023.

\bibitem{yi2025}
P.~Yi, W.~Cheng, J.~Wang, J.~Pan, Y.~Ouyang, and W.~Zhang, ``Intelligent multi-link EDCA optimization for delay-bounded QoS in Wi-Fi 7,'' arXiv:2509.25855, 2025.

\bibitem{zuo2025}
Z.~Zuo, D.~Wang, X.~Nie, X.~Pan, M.~Deng, and M.~Ma, ``PDCF-DRL: a contention window backoff scheme based on deep reinforcement learning for differentiating access categories,'' \emph{J.~Supercomput.}, vol.~81, art.~213, 2025.

\bibitem{du2024}
X.~Du, X.~Fang, R.~He, L.~Yan, L.~Lu, and C.~Luo, ``Federated deep reinforcement learning-based intelligent channel access in dense Wi-Fi deployments,'' arXiv:2409.01004, 2024.

\bibitem{li2024}
Q.~Li, B.~Lv, Y.~Hong, and R.~Wang, ``ReinWiFi: application-layer QoS optimization of WiFi networks with reinforcement learning,'' arXiv:2405.03526, 2024.

\bibitem{jain1984}
R.~Jain, D.~Chiu, and W.~Hawe, ``A quantitative measure of fairness and discrimination for resource allocation in shared computer systems,'' DEC Research Report TR-301, 1984.

\bibitem{ieee80211ax}
IEEE Std 802.11ax-2021, ``Enhancements for High-Efficiency WLAN,'' IEEE, 2021.

\end{thebibliography}

\end{document}
